\input{../tex-inputs/latex-leseansicht-vorspann}

\section[Arthur Schnitzler an Gustav Schwarzkopf, 27. 2. 1903]{L04442 Arthur Schnitzler an Gustav Schwarzkopf, 27. 2. 1903}
\nopagebreak\mylabel{L04442v}
\rehead{ }\normalsize\beginnumbering\briefempfaengerindex{Schwarzkopf, Gustav@\textsc{Schwarzkopf, Gustav}!zzzSchnitzler, Arthur@\emph{von Arthur Schnitzler}!1903-02-272@{27. 2. 1903}|(be}
\toendnotes[C]{\smallbreak\pagebreak[2]}
\correspDesc{Versand  durch Arthur Schnitzler am 27. 2. 1903 in Berlin
\newline{}Übermittlung  am 27. 2. 1903 in Berlin
\newline{}Zustellung  am 28. 2. 1903 in I., Innere Stadt
\newline{}Erhalt  durch Gustav Schwarzkopf im Zeitraum [28. 2. 1903 – 4. 3. 1903?] in Wien}\toendnotes[C]{\smallbreak}
\Standort{DLA, A:Schnitzler, HS.1985.1.1897.}
\physDesc{Bildpostkarte, 84 Zeichen
\newline{}Handschrift: Bleistift, deutsche Kurrent
\newline{}Versand: 1) Stempel: »\nobreak{}\oindex{Berlin@\textbf{Berlin}|pwk}B\textcolor{gray}{e}{[}rli{]}n, 27. 2. 03, 6–7 N\nobreak{}«.   2) Stempel: »\nobreak{}\oindex{I., Innere Stadt@\textbf{I., Innere Stadt}|pwk}Wien 1\textcolor{gray}{/1 1}, 28 2 03, \textcolor{gray}{×}\-\textcolor{gray}{×}\-\textcolor{gray}{×} N, \textcolor{gray}{bestellt}\nobreak{}«. }\toendnotes[C]{\smallbreak}\pstart{}{\pb}\textsc{Herrn} Guſtav Schwarzkopf\pend{}\pstart{}Wien\oindex{Wien@\textbf{Wien}|pw}\pend{}\pstart{}I. Tiefer Graben 23\oindex{Wien@\textbf{Wien}!I., Innere Stadt@\textbf{I., Innere Stadt}!Tiefer Graben 23@\textbf{Tiefer Graben 23}|pw}.\pend{}{\bigskip}
\pstart
           \noindent{}\raggedleft{}{\pb}\textcolor{gray}{\textbf{»Am Potsdamer Platz\oindex{Potsdamer Platz@\textbf{Potsdamer Platz}|pw}.{[}«{]} Dienstmann. Droschkenkutscher.}}\pend
           \vspace{1em}
\pstart
           \noindent{}{\pb}Herzl Gruſs. Ihr \spacefill\mbox{A. S.}\pend
           
\pstart
           \noindent{}\label{T_L04442-1v}\edtext{+ Hier wohn ich.}{\lemma{\textnormal{\emph{+ Hier wohn ich.}}}\Cendnote{\textnormal{Schnitzler markiert das Palasthotel\oindex{Palasthotel Berlin@\textbf{Palasthotel Berlin}|pwk} auf der Fotografie mit einem Kreuz.}}}\label{T_L04442-1}\pend
           \selectlanguage{ngerman}\endnumbering\briefempfaengerindex{Schwarzkopf, Gustav@\textsc{Schwarzkopf, Gustav}!zzzSchnitzler, Arthur@\emph{von Arthur Schnitzler}!1903-02-272@{27. 2. 1903}|)be}\mylabel{L04442h}
\begin{anhang}
\end{anhang}\newcommand{\dateiname}{L04442}\newcommand{\titel}{Arthur Schnitzler an Gustav Schwarzkopf, 27. 2. 1903}\newcommand{\editorInnen}{Herausgegeben von Jahnke, SelmaMüller, Martin Anton}\input{../tex-inputs/latex-leseansicht-abspann}
